\documentclass{beamer}
\usepackage[utf8]{inputenc}
\usepackage[T1]{fontenc}

\usepackage{hyperref}
\usepackage{csquotes}
\usepackage[
	backend=biber,
	style=numeric,
	citestyle=numeric,
	sorting=none,
	url=false]{biblatex}

\setbeamertemplate{bibliography item}[text]
\renewcommand*{\bibfont}{\scriptsize}


\usepackage{array}
\usepackage{tabularx}
\usepackage{lipsum}


\setcounter{tocdepth}{3}

\usepackage{bbold}
\usepackage{scalerel}[2016/12/29]

\newcolumntype{M}[1]{>{$\displaystyle\qquad}p{#1}<{$}}

\usepackage{appendixnumberbeamer}

\usepackage{graphicx}
\graphicspath{{Im/}}


\newcounter{FootlineChoice} 
% Counter to allow the user to choose between a few footlines implemented by owner

\newcounter{DepartmentLogo}
% Counter to allow the user to choose whether to display the Department Logo in NE corner

\newcounter{FacultyLogo}
% Counter to allow the user to choose whether to display the Faculty Logo in SE corner

\usetheme{ULB}
\usetikzlibrary{positioning, arrows.meta, fit, backgrounds}

\setcounter{FootlineChoice}{1}


\title{Advenced Databases\\Object-oriented databases with ObjectDB
and ObjectBox}
\date{14 December 2026}
\author[]{
    Amine Ayadi\texorpdfstring{\\}{, }
    Martin Gouverneur\texorpdfstring{\\}{, }
    Filipe Pereira Martins\texorpdfstring{\\}{, }
    Ethan Van Ruyskensvelde
}
\institute[]
{
    Université Libre de Bruxelles
}


\begin{document}

\begin{frame}[plain,noframenumbering]
    \titlepage
\end{frame}

\begin{frame}{Introduction}
    \begin{block}{Context \& Motivation}
        Modern software is built around \textbf{object-oriented models}, yet persistent storage still relies mostly on \textbf{relational table}, forcing continuous translation between two paradigms.
    \end{block}

    \vspace{0.5em}
    \begin{itemize}
        \item \textbf{Project Goal:} Evaluate whether \textbf{Object-Oriented Databases (OODBMS)} are the technology of choice for graph-heavy applications compared to a relational baseline (\textbf{PostgreSQL}).
        \item     \textbf{Studied OODBMS:} \textbf{ObjectDB} and \textbf{ObjectBox}.

    \end{itemize}
\end{frame}

\begin{frame}
    \vfill
    \begin{center}
        {\usebeamercolor[fg]{title}\LARGE \textbf{Introduction to Object-Oriented Databases}}\\[0.8em]
    \end{center}
    \vfill
\end{frame}

% ====================================================================
% Foundations of Object-Oriented Databases (~5 min)
% ====================================================================

\begin{frame}{Origins and motivation}
    Applications are written with \textbf{object-oriented languages}, but data is stored in \textbf{relational tables}
    $\Rightarrow$ the \textbf{object-relational impedance mismatch}.

    \vspace{0.3em}
    \begin{columns}[T]
        \begin{column}{0.48\textwidth}
            \textbf{Object}
            \begin{itemize}\small
                \item Has its own \textbf{identity}, separate from its values
                \item Two identical objects are still two objects
                \item Links are \textbf{pointers}: follow them in one step
            \end{itemize}
        \end{column}
        \begin{column}{0.48\textwidth}
            \textbf{Row}
            \begin{itemize}\small
                \item Identified by its \textbf{primary key}, which is part of the data
                \item Change the key $\rightarrow$ it's a different row
                \item Links are \textbf{foreign keys}, rebuilt by \textbf{joins}
            \end{itemize}
        \end{column}
    \end{columns}

    \vspace{-0.2em}
    \begin{center}
        \begin{tikzpicture}[
                font=\scriptsize,
                obj/.style={draw, rounded corners, align=center, inner sep=2pt,
                minimum width=1.5cm, minimum height=0.5cm, fill=blue_hl_bg_1!40},
                ptr/.style={-{Stealth[length=1.6mm]}, thick, blue_hl_bg_2},
                match/.style={-{Stealth[length=1.6mm]}, dashed, gray!70!black},
                tbl/.style={draw, inner sep=2pt, font=\scriptsize\ttfamily, minimum width=3.6cm}
            ]
            % object graph
            \node[obj] (cust)  at (0,0)      {\texttt{Customer}};
            \node[obj] (order) at (0,-0.9)   {\texttt{Order}};
            \node[obj] (l1)    at (-0.9,-1.8){\texttt{Line}};
            \node[obj] (l2)    at ( 0.9,-1.8){\texttt{Line}};
            \draw[ptr] (order) -- (cust);
            \draw[ptr] (order) -- (l1);
            \draw[ptr] (order) -- (l2);
            \node[font=\scriptsize\itshape] at (0,-2.45) {links are references};

            % tables
            \node[tbl] (tcust) at (5.2,0)    {customers(id, name)};
            \node[tbl] (tord)  at (5.2,-0.9) {orders(id, customer\_id)};
            \node[tbl] (tline) at (5.2,-1.8) {lines(id, order\_id)};
            \draw[match] (tord.east)  to[bend right=40] node[right] {join} (tcust.east);
            \draw[match] (tline.east) to[bend right=40] node[right] {join} (tord.east);
            \node[font=\scriptsize\itshape] at (5.2,-2.45) {links are matching values};
        \end{tikzpicture}
    \end{center}
\end{frame}

\begin{frame}{Origins and motivation: the translation layer}
    \textbf{Other OO features with no relational equivalent:}
    \begin{itemize}\small
        \item \textbf{Inheritance} $\rightarrow$ flattened onto tables (storage space vs.\ query speed)
        \item \textbf{Collections} $\rightarrow$ moved to separate tables (one value per column)
    \end{itemize}

    \vspace{0.3em}
    \begin{center}
        \begin{tikzpicture}[
                font=\scriptsize,
                box/.style={draw, rounded corners, align=center, minimum height=0.75cm,
                minimum width=2.1cm, inner sep=3pt},
                app/.style={box, fill=blue_hl_bg_1!50},
                map/.style={box, fill=red!12, dashed},
                db/.style={box, fill=gray!15},
                arr/.style={{Stealth[length=1.6mm]}-{Stealth[length=1.6mm]}, thick}
            ]
            \node[app] (a1) at (0,0)   {Application\\objects};
            \node[map] (m1) at (3,0)   {Hand-written\\mapping / ORM};
            \node[db]  (d1) at (6,0)   {Relational\\tables};
            \draw[arr] (a1) -- (m1);
            \draw[arr] (m1) -- (d1);

            \node[app] (a2) at (0,-1.2) {Application\\objects};
            \node[db]  (d2) at (6,-1.2) {Object\\database};
            \draw[arr, blue_hl_bg_2] (a2) -- node[above, font=\scriptsize\itshape] {stored as they are} (d2);
        \end{tikzpicture}
    \end{center}

    \vspace{0.2em}
    \begin{block}{The mapping layer}
        \small
        Belongs neither to the application nor to the database, must be kept in step with the schema,
        and brings its own problems: \textbf{loading too much data}, or \textbf{one query per collection element}.
    \end{block}

    \begin{center}
        \textbf{OODBMS: store objects directly $\Rightarrow$ the middle layer disappears.}
    \end{center}
\end{frame}

\begin{frame}{Fundamental Principles of OODBMS (1/2)}
    \small
    \begin{columns}[T]
        \begin{column}{0.48\textwidth}
            \begin{block}{Object identity (OID)}
                Assigned by the database, unique, never changes.
                Independent of the values: what a primary key cannot provide.
            \end{block}
            \begin{block}{Inheritance}
                The class hierarchy is part of the stored schema, no mapping strategy needed
                (not used in our model).
            \end{block}
        \end{column}
        \begin{column}{0.48\textwidth}
            \begin{block}{Complex objects}
                Attributes can hold other objects, lists, sets or maps, nested arbitrarily.
                The structure is stored as the application sees it.
            \end{block}
            \begin{block}{References}
                Relationships are stored as references. Following one is a direct jump:
                \textbf{no join needed}.
            \end{block}
        \end{column}
    \end{columns}
\end{frame}

\begin{frame}{Fundamental Principles of OODBMS (2/2)}
    \small
    \begin{columns}[c]
        \begin{column}{0.52\textwidth}
            \begin{block}{Persistence by reachability}
                Saving an object saves everything it can reach, recursively.
                A whole structure is stored in \textbf{one operation}.
            \end{block}
        \end{column}
        \begin{column}{0.44\textwidth}
            \centering
            \begin{tikzpicture}[
                    font=\scriptsize,
                    obj/.style={draw, rounded corners, align=center, inner sep=2pt,
                    minimum width=1.4cm, minimum height=0.5cm, fill=blue_hl_bg_1!40},
                    root/.style={obj, fill=blue_hl_bg_2, text=white, very thick},
                    ptr/.style={-{Stealth[length=1.6mm]}, thick}
                ]
                \node[root] (order) at (0,0)       {\texttt{Order}};
                \node[obj]  (cust)  at (-1,-1)     {\texttt{Customer}};
                \node[obj]  (line)  at (1,-1)      {\texttt{Line}};
                \node[obj]  (prod)  at (1,-1.9)    {\texttt{Product}};
                \draw[ptr] (order) -- (cust);
                \draw[ptr] (order) -- (line);
                \draw[ptr] (line)  -- (prod);
                \node[anchor=west, font=\scriptsize\itshape] at (0.75,0) {\texttt{persist()}};
                \begin{scope}[on background layer]
                    \node[draw, dashed, rounded corners, inner sep=5pt,
                    fit=(cust)(line)(prod)] (reach) {};
                \end{scope}
                \node[anchor=north, font=\scriptsize\itshape] at (reach.south) {stored too};
            \end{tikzpicture}
        \end{column}
    \end{columns}

    \vspace{0.1em}
    \begin{block}{Navigation \emph{and} declarative queries}
        \textbf{Navigation:} start from an object, follow its references.\\
        \textbf{Queries:} describe a condition, the database finds the matches.\\
        Relational databases only offer queries; OODBMS offer both.
    \end{block}

    \begin{block}{Transparent persistence}
        A loaded object is an ordinary language object, ready to use. Modifying it inside a
        transaction is enough to store the change: \textbf{no translation layer}.
    \end{block}
\end{frame}

\begin{frame}{Standards and norms}
    \begin{center}
        \begin{tikzpicture}[
                font=\scriptsize,
                dot/.style={circle, fill=blue_hl_bg_2, inner sep=2.2pt},
                lbl/.style={align=center, text width=1.6cm, anchor=north}
            ]
            \draw[thick, -{Stealth[length=2mm]}, blue_hl_bg_2] (-0.4,0) -- (8.8,0);
            \foreach \x/\y/\lab in {0/1989/{OODB\\Manifesto},
                                  2.1/1991/{ODMG\\founded},
                                  4.2/2000/{ODMG 3.0\\(ODL, OQL)},
                                  6.3/2001/{Java binding\\$\rightarrow$ JDO},
                                  8.4/2006/{JPA\\succeeds JDO}} {
                \node[dot] at (\x,0) {};
                \node[above=3pt, font=\scriptsize\bfseries] at (\x,0) {\y};
                \node[lbl] at (\x,-0.2) {\lab};
            }
        \end{tikzpicture}
    \end{center}

    \vspace{0.6em}
    \begin{itemize}
        \item The \textbf{Manifesto} defines what an OODB is: object features + usual database features
        \item \textbf{ODMG} tries to standardise it\ldots{} then disbands, its Java part ends up in JPA
    \end{itemize}

    \vspace{0.6em}
    \begin{alertblock}{Main point: no common standard today}
        \textbf{ObjectDB} uses JPA (an API, not a storage model), \textbf{ObjectBox} has its own API.\\
        $\Rightarrow$ Code written for one cannot be moved to the other.
    \end{alertblock}
\end{frame}

\begin{frame}{Relational Model vs. Object-Oriented Model}
    \small
    \renewcommand{\arraystretch}{1.3}
    \begin{center}
        \begin{tabularx}{\textwidth}{|>{\bfseries}l|>{\raggedright\arraybackslash}X|>{\raggedright\arraybackslash}X|}
            \hline
            & \textbf{Relational} & \textbf{Object-oriented} \\
            \hline
            Data model       & tables of rows, atomic values  & objects with attributes \& collections \\
            \hline
            Identity         & primary key from the data      & identifier assigned by the system \\
            \hline
            Relationships    & foreign keys, rebuilt by joins & references, followed in one step \\
            \hline
            Inheritance      & mapped onto tables             & part of the stored schema \\
            \hline
            Query language   & SQL, standard \& universal     & specific to each product \\
            \hline
            Schema evolution & mature tooling, migrations     & depends on the product \\
            \hline
            Ecosystem        & very large                     & small \\
            \hline
        \end{tabularx}
    \end{center}
\end{frame}

\begin{frame}
    \vfill
    \begin{center}
        {\usebeamercolor[fg]{title}\LARGE \textbf{Introduction to ObjectDB and ObjectBox}}\\[0.8em]
    \end{center}
    \vfill
\end{frame}

\begin{frame}{ObjectDB}
\end{frame}

\begin{frame}{ObjectBox}
\end{frame}

\begin{frame}{PostgreSQL (Baseline)}
\end{frame}

\begin{frame}
    \vfill
    \begin{center}
        {\usebeamercolor[fg]{title}\LARGE \textbf{Application Benchmark \\ \& \\
        \vspace{0.2em} Comparative Analysis}}\\[0.8em]
    \end{center}
    \vfill
\end{frame}

\begin{frame}{Application domain}
    \textbf{Context:} Role-playing game (RPG) modeled after the \textit{Pokémon} game.
    \vspace{0.4em}

    \begin{block}{Two Distinct Categories of Data}
        \begin{itemize}
            \item \textbf{Static reference data:} Immutable Pokémon entities (name, type, speed\dots) and item definitions (name, description).
            \item \textbf{Dynamic player state:} Mutable account credentials, active and customized teams, and inventory.
        \end{itemize}
    \end{block}

    \vspace{0.4em}
    \textbf{Why an ideal case study (OODBMS vs. RDBMS)?}
    \begin{itemize}
        \item \textbf{Structural complexity:} Nested variable-length collections (team lists, item-quantity associations), object composition, and complex entity references.
        \item \textbf{Mixed workloads:} Single-profile reads/updates alongside cross-player global searches (e.g., matchmaking inspecting all teams for legendary Pokémon).
    \end{itemize}
\end{frame}

\begin{frame}{Conceptual schema and data models}
\end{frame}

\begin{frame}{Software architecture}
\end{frame}

\begin{frame}{Dataset generation}
    \small
    \begin{block}{\normalsize Authentic Reference Datasets (\textit{Kaggle})}
        \begin{itemize}
            \item \texttt{\href{https://www.kaggle.com/datasets/mlomuscio/pokemon}{pokemons.csv}}: 800 Pokémon with full stats (\texttt{HP}, \texttt{Attack}, \texttt{Defense}, \texttt{Sp. Atk}, \texttt{Sp. Def}, \texttt{Speed}, \texttt{Generation}, \texttt{Legendary}).
            \item \texttt{\href{https://www.kaggle.com/datasets/mylesoneill/pokemon-sun-and-moon-gen-7-stats?select=items.csv}{items.csv}}: 920 game items with names and descriptions.
        \end{itemize}
    \end{block}

    \begin{block}{\normalsize Deterministic Workload Generator (\texttt{DatasetGenerator.java})}
        \textbf{Scale:} $1\text{K}$, $10\text{K}$, $100\text{K}$, and $1\text{M}$ players \hfill \textbf{Reproducibility:} Fixed \texttt{SEED = 12345L}
        \begin{itemize}
            \item \textbf{Identity:} Unique \texttt{Player\_i} ($i \in \{1, \dots, N\}$) + deterministic \texttt{password\_hash}.
            \item \textbf{Collection:} $1$ to $10$ distinct Pokémon (sampled without replacement).
            \item \textbf{Teams:} $1$ to $10$ custom teams ($1$–$6$ slots from owned Pokémon).
            \item \textbf{Inventory:} $1$ to $10$ distinct items (quantities from $1$ to $10$).
        \end{itemize}
    \end{block}
\end{frame}

\begin{frame}{Experimentation and benchmarking protocol}
\end{frame}

\begin{frame}{Result and comparative analysis}
\end{frame}

\begin{frame}{Overall Assessment: Pro and Cons}
\end{frame}

\begin{frame}{Conclusion}
    \begin{itemize}
        \item Eliminating object-relational impedance mismatch yields
            substantial performance gains for graph-oriented workloads.
        \item \textbf{ObjectDB:} Best write throughput.
        \item \textbf{ObjectBox:} Fastest point reads.
        \item \textbf{PostgreSQL:} Essential when shared SQL access or $>1\text{M}$ players without licensing costs is required.
    \end{itemize}
    \vspace{1em}
    \begin{center}
        \textbf{Object-oriented databases prove to be the technology of
        choice for our RPG application domain.}
    \end{center}
\end{frame}

\end{document}
